\chapter{Tổng quan giải pháp}
\label{chap:solution-overview}

\section{Hệ thống giải quyết việc gì?}

SmartWater kết nối lớp thiết bị đo với đội vận hành. Hệ thống thu nhận TDS và cảnh báo từ MQTT, lưu lịch sử telemetry để theo dõi theo MCU/thiết bị, đồng thời cung cấp các chức năng quản lý tài sản, khách hàng và bảo trì trong SmartWater Admin. Người dùng không cần làm việc trực tiếp với broker hoặc database để xem dữ liệu vận hành.

\begin{table}[htbp]
\centering
\small
\begin{tabularx}{\textwidth}{p{0.22\textwidth}X p{0.16\textwidth}}
\toprule
\textbf{Nhóm chức năng} & \textbf{Kết quả người dùng nhận được} & \textbf{Bằng chứng} \\
\midrule
Telemetry & Nhận, chuẩn hóa, lưu và tra cứu TDS/cảnh báo theo \texttt{mcu\_id}. & E-SIM-03, E-DM-02--E-DM-04, E-DB-02 \\
Giám sát & Chọn MCU, lọc lịch sử và xem chuỗi TDS theo thời gian. & E-DM-04, E-DM-06 \\
Quản trị & Quản lý thiết bị, sản phẩm, kho, khách hàng, hợp đồng, nhân viên và bảo trì. & E-ADM-04 \\
Khai thác tại Admin & Xem telemetry và cảnh báo trong ngữ cảnh thiết bị. & E-ADM-05 \\
\bottomrule
\end{tabularx}
\caption{Chức năng được mô tả theo kết quả sử dụng}
\label{tab:solution-function-outcomes}
\end{table}

\section{Luồng telemetry}

Hệ thống SmartWater được triển khai theo nhiều ranh giới chức năng, giao tiếp thông qua
MQTT, HTTP và MySQL. Luồng xử lý telemetry được tổ chức như sau:

\begin{enumerate}
    \item Simulator ESP32 tạo và gửi dữ liệu telemetry qua MQTT.
    \item HiveMQ Cloud đóng vai trò MQTT Broker.
    \item Device Monitor nhận dữ liệu qua kết nối MQTT over WebSocket (WSS).
    \item Slim API xác thực và lưu dữ liệu vào bảng \texttt{telemetry}.
    \item SmartWater Admin truy vấn dữ liệu để hiển thị giá trị TDS và trạng thái cảnh báo trên trang chi tiết thiết bị.
\end{enumerate}

Mã evidence của luồng này là E-SIM-03, E-DM-02--E-DM-06, E-DB-02 và E-ADM-05. Đường dẫn source đầy đủ được tra tại Phụ lục~\ref{app:source-evidence}.
\begin{landscape}
\begin{figure}[p]
\centering
\resizebox{0.98\linewidth}{!}{%
\begin{tikzpicture}[
  node distance=8mm,
  box/.style={draw=SmartWaterBlue,rounded corners=3pt,align=center,text width=3.15cm,minimum height=1.65cm,fill=SmartWaterBlue!7,font=\small},
  external/.style={box,draw=StatusPartial,dashed,fill=StatusPartial!10},
  partial/.style={box,draw=StatusPartial,fill=StatusPartial!8},
  legacy/.style={box,draw=StatusPlanned,dashed,fill=black!6},
  flow/.style={-{Latex[length=2.5mm]},thick,draw=SmartWaterNavy},
  branch/.style={flow,dashed,draw=StatusPartial}
]
\node[partial] (sim) {Simulator ESP32\\\texttt{mcu\_id, tds, alert}\\Partially Implemented};
\node[external,right=of sim] (broker) {HiveMQ\\Broker bên ngoài\\Configured};
\node[box,right=of broker] (browser) {Device Monitor\\mqtt.js / WSS\\Implemented};
\node[box,right=of browser] (slim) {Slim API\\Normalize + persist\\Implemented};
\node[box,right=of slim] (db) {MySQL canonical\\\texttt{telemetry.mcu\_id}\\Implemented};
\node[partial,right=of db] (admin) {SmartWater Admin\\Hiển thị theo thiết bị\\Partially Implemented};

\node[legacy,below=24mm of sim] (firmware) {Firmware chính\\Payload \texttt{device\_id}\\Contract legacy};
\node[legacy,below=24mm of browser] (worker) {.NET MQTT Worker\\DTO \texttt{device\_id}\\Partially Implemented};
\node[legacy,right=12mm of worker] (legacydb) {MySQL legacy\\\texttt{devices, device\_data}\\Không canonical};

\draw[flow] (sim) -- node[above]{TLS} (broker);
\draw[flow] (broker) -- node[above]{WSS} (browser);
\draw[flow] (browser) -- node[above]{HTTP} (slim);
\draw[flow] (slim) -- node[above]{PDO} (db);
\draw[flow] (db) -- node[above]{ORM} (admin);
\draw[branch] (firmware.east) -| node[pos=0.22,below]{mismatch} (broker.south);
\draw[branch] (broker.south) |- node[pos=0.55,above]{SUB} (worker.north);
\draw[branch] (worker) -- node[above]{Dapper} (legacydb);
\end{tikzpicture}%
}

\vspace{3mm}
\small\textbf{Chú giải:} nét liền là luồng canonical theo source; nét đứt là nhánh
legacy hoặc ranh giới chưa tương thích. Broker là phụ thuộc bên ngoài repository.
\caption{Kiến trúc tổng thể và các ranh giới hợp đồng SmartWater}
\label{fig:system-architecture}
\end{figure}
\end{landscape}


\section{Mục tiêu và năng lực đã có}

Giải pháp hiện có hỗ trợ ba mục tiêu có bằng chứng source. Chi tiết đường dẫn được tập trung tại Phụ lục~\ref{app:source-evidence}, để phần này giữ trọng tâm ở năng lực hệ thống:

\begin{itemize}
  \item thu nhận TDS và trạng thái cảnh báo từ publisher MQTT; simulator tạo payload \texttt{mcu\_id} và \texttt{telemetry.tds}, còn firmware chính phát thêm cảnh báo trong object telemetry (E-SIM-03, E-FW-04);
  \item giám sát gần thời gian thực và lưu lịch sử qua dashboard/API Device Monitor; các route đọc, tổng hợp và ghi telemetry được khai báo trong \path|Project/Device-monitor/public/index.php| (E-DM-02);
  \item khai thác lịch sử telemetry theo thiết bị trong Admin; \texttt{DeviceController::show} truy vấn quan hệ \texttt{Device::telemetry} và view \path|resources/views/devices/show.blade.php| hiển thị đồ thị TDS/cảnh báo.
\end{itemize}

Hệ thống đã được live test thành công theo chuỗi telemetry MCU $\rightarrow$ HiveMQ $\rightarrow$ Device Monitor $\rightarrow$ MySQL. MCU publish dữ liệu lên HiveMQ, Device Monitor nhận message, chuẩn hóa payload và lưu bản ghi telemetry vào database; vì vậy sơ đồ thể hiện cả cấu trúc source và luồng tích hợp đã được xác nhận trong môi trường chạy thực tế.

\section{Phạm vi và trạng thái}

\begin{table}[htbp]
\centering
\small
\begin{tabular}{p{0.25\textwidth}p{0.19\textwidth}p{0.47\textwidth}}
\hline
\textbf{Khối} & \textbf{Trạng thái} & \textbf{Căn cứ} \\
\hline
Simulator telemetry & Implemented & Có publisher định kỳ và payload \texttt{mcu\_id} dạng chuỗi cùng \texttt{telemetry.tds} theo contract canonical (E-SIM-01--E-SIM-04). \\
HiveMQ broker & Configured & Có cấu hình và client kết nối nhưng không tìm thấy manifest server-side trong repository (E-MQ-06). \\
Device Monitor & Implemented & mqtt.js subscribe qua WSS, Slim API chuẩn hóa và PDO repository ghi/đọc \texttt{telemetry} (E-DM-02--E-DM-06). \\
SmartWater Admin & Partially Implemented & Các module lõi, API telemetry và quan hệ MCU có source; tổng thể còn view thiếu, mock data và mismatch validation/navigation (E-ADM-02--E-ADM-07). \\
Database dùng chung & Implemented & Migration chuẩn định nghĩa \texttt{telemetry.mcu\_id} dạng chuỗi, có index riêng và index kết hợp với timestamp (E-DB-01--E-DB-03). \\
Firmware chính & Implemented & Có sensor/task/MQTT/command; telemetry dùng \texttt{mcu\_id} dạng chuỗi và object \texttt{telemetry} chứa TDS/cảnh báo trên topic canonical (E-FW-03, E-FW-04). \\
.NET MQTT worker & Partially Implemented & Subscribe và lưu MySQL đã có, nhưng DTO/schema riêng dùng \texttt{device\_id}, \texttt{devices}, \texttt{device\_data} (E-MQ-01--E-MQ-05). \\
\hline
\end{tabular}
\caption{Phạm vi và trạng thái giải pháp theo bằng chứng source}
\label{tab:solution-scope-status}
\end{table}

\section{Quyết định kiến trúc chi phối}

Migration dùng chung là nguồn chuẩn của cấu trúc dữ liệu. Theo \texttt{2026\_07\_13\_000000\_create\_telemetry\_table.php}, \texttt{mcu\_id} là chuỗi tối đa 50 ký tự, có index và không được migration telemetry ràng buộc khóa ngoại; Device Monitor yêu cầu giá trị chuỗi không rỗng cùng giới hạn độ dài (E-DB-02, E-DM-04).

Live test đã sử dụng định danh MCU hợp lệ và xác nhận phép ghi telemetry thành công vào database. \texttt{SimulatorApp::publishTelemetry} gửi trực tiếp \texttt{SimulatorConfig::MCU\_ID} dạng chuỗi, phù hợp với repository hiện tại.

SmartWater Admin cung cấp \path|POST /api/telemetry| như một cổng ingest độc lập, nhưng source Device Monitor đang POST vào chính route Slim \path|/api/telemetry|; không có lời gọi từ Device Monitor sang API Laravel. Hai cổng cùng hướng đến bảng chuẩn nhưng là hai đường HTTP khác nhau (E-DM-02, E-DM-06, E-ADM-05).
